\section{What the Experiments Actually Establish}
EEGraft brings together an implemented decoder and evidence about how EEG classification should be assessed. The custom model combines a frozen pretrained representation with learnable temporal features; its source and external comparisons reveal different observed model rankings. A separate five-minute SVM analysis and the Isfahan negative-control experiment quantify persistent incorrect decisions. The external CAUEEG investigation shows that the same frozen XGBoost scores can give very different eye-condition consistency under two decision thresholds. These results support a useful distinction: discrimination, calibration, temporal consistency and correctness are related, but none can replace the others.

The scientific value is not a single winning percentage. It is the documented connection between an engineering design, a strong classical comparator, an independently held-out source test, frozen transfer, and checks for hidden error patterns. The research questions can therefore be answered with direct observations rather than relying on model complexity as evidence of quality.

\section{Why a Decoder Is Still a Technical Contribution}
CBraMod itself was introduced by Wang et al. \cite{wang2025cbramod}; hybrid convolutions, attention and residual fusion also have earlier precedents \cite{dice2023,song2023}. EEGraft's engineering work is the particular trainable decoder that connects frozen CBraMod tokens to a two-scale raw-signal pathway, a sigmoid-controlled additive residual, and class-specific attention over 76 tokens. The implemented modules, training configuration and ablations make that design reproducible and directly testable.

The 52-person development comparison provides evidence about this configuration: full DSRF obtained 77.06\% balanced accuracy, compared with 73.16\% for a capacity-matched gated blend and lower values for the single-path and no-skip variants. These are informative within-protocol differences. They motivate replication with other seeds and cohorts rather than establishing a population-wide architecture ranking. A documented and ablated model remains a substantive engineering contribution even when it does not win the separate locked source comparison.

\section{Why the Classical Baseline Is Important}
XGBoost reached nearly the same development BAcc as DSRF and exceeded it on the locked 13-person NEMAR test. It is therefore a meaningful comparator, not a straw-man baseline. On the external CAUEEG test, however, its original threshold 0.50 predicted broad Dementia for 51 of 57 recordings and identified only six of 35 Normal recordings. Reapplying the training-selected 0.70 threshold improved observed balanced accuracy and increased eye-state label changes. This illustrates why thresholding should not be treated as an invisible part of the architecture. Even a high percentage of stable decisions may reflect the decision rule or score distribution rather than trustworthy disease discrimination.

\section{Temporal Stability Is Not Correctness}
The negative-control Isfahan results and the CAUEEG unanimous-window counts make the same conceptual point in different settings. In Isfahan EEG\_1, a high stable fraction coexisted with below-chance first-minute balanced accuracy. In CAUEEG, XGBoost was unanimous and wrong in 12 recordings at the primary rule; DSRF had fewer unanimous records, but greater aggregate accuracy. Neither result makes XGBoost inherently unsafe nor DSRF inherently clinically safer. Temporal or condition stability depends on the comparison interval, the number of windows, the aggregation method and the threshold. Any claim that one model is reliably superior would require predefined metrics and larger independent cohorts.

\section{External Validity and Diagnostic-Label Shift}
NEMAR is an AD/CN research case-control dataset. The CAUEEG Dementia task covers broader diagnostic labels assigned in a different clinical context. Even though both have 19 compatible EEG inputs, differences in acquisition reference, age, conditions, event sequences and subject mix remain. Transferring frozen models without target weight updates tests performance under these combined shifts; it does not separate their causal contributions or validate the original AD-specific decision in a new hospital population. The official CAUEEG overlap-reduced protocol removes known cross-partition reuse, but full de-identified patient linkage is not independently available for confidence intervals on \emph{people}. We therefore report recording-level outcomes only \cite{kim2023caueeg,caueegofficial}.

The external frozen-test confidence interval is wide. The primary DSRF-minus-XGBoost BAcc interval ranges from $-3.31$ to $+22.08$ percentage points, crossing zero. Its observed test difference cannot be treated as established superiority. Nor can published CAUEEG systems trained on CAUEEG labels be directly ranked against models frozen from NEMAR without accounting for their fundamentally different tasks and training conditions.

\section{Age, Event Conditions and Calibration Confounding}
The CAUEEG age audit raises a particularly strong validity concern: the Normal and Dementia recordings differ substantially in mean age and an age-only classifier already reaches validation AUROC 0.8397. Adding DSRF increases validation AUROC by just 0.0118, with a 95\% interval that includes negative and positive gains. It is therefore plausible that EEG scores add information, but an independently verified age-adjusted improvement has not been demonstrated. Some of the broader Dementia labels also represent diagnoses other than AD, further limiting a disease-specific conclusion.

Eye-open and eye-closed scores differ on eligible recordings, but recording time and eye state are correlated. Training associations and within-recording linear adjustments do not eliminate hidden confounding. In the smaller validation set the adjusted coefficient intervals cross zero. The findings motivate controlled acquisition-condition experiments rather than justify causal statements about eye closure or photic stimulation. Calibration reduced reported CAUEEG binary Brier scores but was evaluated post-hoc on a small set and cannot validate risk predictions or patient-level probabilities.

\section{Statistical and Reproducibility Limitations}
The original source locked test includes only seven AD and six control people, so single errors cause large percentage changes. Study A's 230/205 participant-repeat records contain repeated measurements and must not be treated as independent participants. Isfahan release identities were not fully reconciled. CAUEEG's 57 binary test units are recordings with incomplete patient linkage; the bootstrap cannot capture uncertainty from unknown repeated patients or future hospitals. Four-second windows sampled from one recording remain dependent. These constraints qualify p-values, confidence intervals and any generalization claim.

The original project source supports exact decoder inputs, architecture, training settings, final model hashes and source holdout labels. The archived Study A duration row discrepancy was resolved as an audit-script model-filtering error, supported by source CSV checks. A full corrected duration audit and end-to-end raw EEG reproduction remain unperformed in the material supplied for this rewrite. The CAUEEG Code 19--25 totals and figures trace to the user's local execution reports, but the restricted raw test signals, checkpoint binaries and individual scores are not in this Overleaf package. The final report reproduces reported \emph{aggregates} and recomputes some arithmetic; it does not claim independent re-execution of those experiments. Dataset permissions prevent including unapproved CAUEEG signal excerpts.

\section{Engineering, Ethical and Clinical Significance}
The practical outcome of EEGraft is an evaluation approach that another research team can apply: compare a custom decoder to a serious classical baseline, keep people separate across source training and testing, preserve the primary operating rule, show incorrect unanimous predictions, and test whether observed performance survives age and acquisition-condition scrutiny. Saved configurations, separated protocols, aggregate results and provenance references make these checks more traceable than an accuracy-only summary. The distinction matters especially for EEG foundation models, where an apparently strong representation need not produce a calibrated or consistent decision in a new cohort \cite{xiong2026,liu2026compass}.

These conclusions concern research evaluation and reproducibility, not clinical diagnosis. The main threats are a 13-person source holdout, recording-level rather than patient-linked external data, broad-Dementia versus AD-specific labels, post-hoc secondary choices, and the sizeable CAUEEG age imbalance. Those are explicit study boundaries, not reasons to discard the measured findings. Future investigators can address them with patient-linked prospective cohorts, matched ages and acquisition conditions, and prespecified operating thresholds and calibration checks.